\section{Detailed Renewal Protocol}

Consider a stipulated service-batch agreement administered by a renewal desk and executed by a shared dispatch gateway. The desk observes the customer's current outside alternative and approved continuation. It can perform an advance tier and issue one of a small number of dispatch tickets. The gateway implements the terminal tier named by that ticket; its interface does not retain the renewal record. A separate billing audit can retain the contract without making that record an input to the gateway. Limiting tickets can price gateway-state certification, transmission, or a restricted instruction set. This is a stipulated stylized mechanism, not a claim that a particular industry already uses this architecture.

In the expected-participation version, the offer communicates the advance tier, the terminal ticket levels, and their probabilities. The customer evaluates expected remaining payment before the ticket is drawn and either renews or exits at that review. The protocol commits to the advertised draw. After renewal the desk performs the advance tier, draws the ticket, and sends only that ticket to the gateway. The realized ticket may be disclosed to the customer. ``Private'' means that the seed is neither a pre-acceptance conditioning variable nor a free future decoder state; it need not mean that realized service is hidden forever. The customer does not receive a new costless exit option after observing an unfavorable draw within this modeled renewal episode.

This version is meaningful only for an ex ante willingness-to-pay or expected-obligation cap and an agreement that permits the stated realization variation. It is not a way to satisfy an ex post enforceable payment ceiling in expectation. A hard billing limit, a service-safety requirement on every draw, or a review after disclosure instead calls for pathwise feasibility, with the deterministic frontier as its optimum. We make no legal assertion about enforceability. The two feasible sets are explicitly different design alternatives.

Repeated customers who may condition a later renewal on past tickets require that history to enter the later public state, cap, or charged memory. The three-date theorem is not an unqualified repeated-game result. Similarly, a contractible seed disclosed before acceptance changes the conditioning information and hence the participation test. It cannot simply be added as a free source of branch-dependent execution. Section~\ref{sec:r37-resources} gives a bounded-realization variant for an operator who permits some variation but not the full expected-participation lottery.
